\documentclass[11pt,a4paper]{article}

\usepackage[utf8]{inputenc}
\usepackage[T1]{fontenc}
\usepackage{lmodern}
\usepackage[a4paper,margin=2.3cm]{geometry}
\usepackage{hyperref}
\usepackage{longtable}
\usepackage{booktabs}
\usepackage{array}
\usepackage{listings}
\usepackage{xcolor}
\usepackage{enumitem}
\usepackage{fancyhdr}

\hypersetup{colorlinks=true,linkcolor=blue,urlcolor=blue}
\definecolor{codegray}{gray}{0.95}
\lstset{basicstyle=\ttfamily\small,backgroundcolor=\color{codegray},frame=single,breaklines=true,columns=fullflexible}
\pagestyle{fancy}
\fancyhead[L]{HeatingSystemController}
\fancyhead[R]{Fejlesztői dokumentáció}
\fancyfoot[C]{\thepage}

\title{HeatingSystemController\\Fejlesztői dokumentáció}
\author{AlfaOmega projekt}
\date{\today}

\begin{document}
\maketitle
\tableofcontents
\newpage

\section{Dokumentum célja}

Ez a dokumentum a HeatingSystemController firmware továbbfejlesztéséhez szükséges technikai információkat foglalja össze. Ismerteti a futási architektúrát, a fűtési rendszer modelljét, a HTTP API-kat, a konfigurációs formátumot, a SPIFFS fájlrendszert, a bináris állapotnapló formátumát, a frontend adat-visszaépítését, a build környezeteket és az integrációs teszteket.

A projekt fő forrásfájljai:

\begin{itemize}
  \item \texttt{src/main.cpp}: inicializálás, fő ciklus, HTTP és WebSocket életciklus.
  \item \texttt{src/web\_api.cpp}: állapot-, konfigurációs-, hőmérséklet- és vezérlő API-k.
  \item \texttt{src/config.h}: konfiguráció betöltése, validálása és feltöltése.
  \item \texttt{src/HeatingSystemElements/}: kazán, puffer, radiátor, szenzor és vezérlő objektumok.
  \item \texttt{src/state\_data\_logger.cpp}: bináris napi állapotnapló és JSON parser API.
  \item \texttt{config/config.json}: fejlesztői konfigurációs példa.
  \item \texttt{data/config.json}: a firmware fájlrendszerébe feltöltendő konfiguráció.
\end{itemize}

\section{Rendszerarchitektúra}

Az alkalmazás Arduino frameworkre épül, PlatformIO-val fordítható ESP8266 és ESP32 célokra. A firmware a SPIFFS fájlrendszert használja webes fájlok, konfiguráció és runtime naplófájlok tárolására.

\begin{lstlisting}
                 +----------------------+
                 |     HTTP clients     |
                 | Postman / frontend   |
                 +----------+-----------+
                            |
                            v
                 +----------------------+
                 | ESP WebServer       |
                 | main.cpp routes     |
                 +----+-------------+---+
                      |             |
             +--------v--+    +-----v----------------+
             | web_api   |    | state_data_logger    |
             | control   |    | binary log + parser  |
             +-----+-----+    +----------+------------+
                   |                     |
                   v                     v
             +-----------+       +-------------+
             | Heating   |       | SPIFFS      |
             | System    |       | config/web  |
             | state      |       | state_*.bin |
             +-----+-----+       +-------------+
                   |
        +----------+-----------+
        | sensors / pumps / IO |
        +----------------------+
\end{lstlisting}

\subsection{Inicializálási sorrend}

A \texttt{setup()} fő lépései:

\begin{enumerate}
  \item soros port és hardver inicializálása;
  \item SPIFFS és \texttt{config.json} betöltése;
  \item \texttt{HeatingSystem} objektum létrehozása;
  \item WiFi és NTP kliens inicializálása;
  \item HTTP route-ok regisztrálása;
  \item szenzorok és fűtési elemek inicializálása;
  \item \texttt{StateDataLogger::begin()} meghívása;
  \item WebSocket és időzítő indítása.
\end{enumerate}

Az \texttt{loop()} minden ciklusban kezeli a hálózatot, frissíti az NTP időt és a Dallas szenzorokat, feldolgozza a fűtési rendszer jelzéseit, majd meghívja a \texttt{stateDataLogger.capture()} függvényt.

\section{Fűtési rendszer modellje}

A \texttt{HeatingSystem} a konfigurációból létrehozza a fűtési elemeket és a közös \texttt{mergedList} listában tartja őket. A fő elemtípusok:

\begin{itemize}
  \item \textbf{Kazan}: kazán állapotgép, visszatérő-védelem és túlmelegedés-védelem.
  \item \textbf{Puffer}: töltés, kisütés, készenlét és telítettségi állapotok.
  \item \textbf{Radiator}: fogyasztói oldal állapotgépe.
  \item \textbf{HeatingController}: a teljes rendszer módja és hőigénye.
\end{itemize}

A szenzorok a \texttt{Sensor::SensorsValue} mapben jelennek meg az egyedi DS18B20 azonosítójukkal. A kimenetek az \texttt{IOMap} mapben vannak, ahol a kulcs az IO szám, az érték tipikusan 0--100 közötti PWM vagy ON/OFF érték.

A modulállapotok a \texttt{HeatingElement::ElementsStateMap} mapben vannak. Gyakori mezők:

\begin{itemize}
  \item \texttt{state}: állapotnév, például \texttt{CHARGE}, \texttt{DISCHARGE}, \texttt{HEATING};
  \item \texttt{HT}: hőátadási érték, floatként naplózva;
  \item \texttt{NH}: hőigény;
  \item \texttt{iA}: aktív állapot;
  \item \texttt{iOHPA}, \texttt{iRPA}: kazánvédelmi állapotok.
\end{itemize}

\section{Konfiguráció}

A konfiguráció JSON fájlban található a SPIFFS gyökerében \texttt{/config.json} néven. A projektben található példa a \texttt{config/config.json} és \texttt{data/config.json} fájlban van.

\subsection{Globális mezők}

\begin{longtable}{p{0.29\textwidth}p{0.62\textwidth}}
\toprule
Mező & Jelentés \\
\midrule
\texttt{wifiSSID} & STA módú WiFi hálózat neve. \\
\texttt{wifiPassword} & STA módú WiFi jelszó; publikus dokumentációban ne szerepeljen valódi érték. \\
\texttt{AP\_SSID} & Access Point mód hálózatneve. \\
\texttt{AP\_Password} & Access Point jelszava. \\
\texttt{WifiMode} & \texttt{STA}, \texttt{AP} vagy \texttt{STA+AP}. \\
\texttt{hostname} & DHCP/mDNS név. \\
\texttt{comType} & \texttt{API} vagy \texttt{WS}; WebSocket esetén a frontend állapotbroadcastot is kaphat. \\
\bottomrule
\end{longtable}

\subsection{Fűtési elemek}

A \texttt{HeatingSystem} objektumai a \texttt{KAZAN}, \texttt{RADIATOR} és \texttt{PUFER} tömbökben vannak. Egy szenzor mezői:

\begin{lstlisting}[language=json]
{
  "model": "DS12B20",
  "position": "body",
  "id": "2814576B00000000059",
  "offset": -1.0,
  "level": 80
}
\end{lstlisting}

A \texttt{position} értéke \texttt{tour}, \texttt{retour} vagy \texttt{body}. Több body szenzor esetén a \texttt{level} súlyozza az átlagot.

A pumpa fontos mezői a \texttt{name}, \texttt{workingMode}, \texttt{IOnumber}, \texttt{minControlSig}, \texttt{maxControlSig} és opcionálisan a \texttt{dischargepump}. PID használatakor a \texttt{pid.enabled}, \texttt{p}, \texttt{i}, \texttt{d}, \texttt{tempSensorId} és \texttt{setPoint} mezők kötelezőek.

\section{HTTP API}

Az API alapértelmezett portja 80. A route-ok regisztrációja a \texttt{src/main.cpp} \texttt{registerHttpRoutes()} függvényében történik.

\subsection{API-kérés limit}

Az API limiter a \texttt{API\_MIN\_REQUEST\_INTERVAL\_MS} build flaggel konfigurálható. Alapértéke 0, ezért jelenleg ki van kapcsolva.

\begin{lstlisting}
-DAPI_MIN_REQUEST_INTERVAL_MS=250
\end{lstlisting}

Bekapcsoláskor a túl gyors kérés HTTP 429 választ és \texttt{Retry-After} fejlécet ad.

\subsection{Állapot lekérése}

\begin{lstlisting}
GET /state
\end{lstlisting}

A válasz egy JSON objektum. Az IO-k és szenzorok lapos kulcsként, a modulállapotok modulnév alatt jelennek meg.

\begin{lstlisting}[language=json]
{
  "1": 0,
  "3": 100,
  "28000B96A00000000029": 55.4,
  "Puffer950l": {
    "state": "CHARGE",
    "HT": "12.3"
  }
}
\end{lstlisting}

\subsection{Hőmérsékletértékek beállítása}

\begin{lstlisting}
POST /setTemp
Content-Type: application/json
\end{lstlisting}

Példa:

\begin{lstlisting}[language=json]
{
  "28000B96A00000000029": 55.0,
  "2814576B00000000059": 48.5
}
\end{lstlisting}

Sikeres válasz: \texttt{"status":"success"}. Ezt a tesztkörnyezet használja az érzékelőértékek szimulálására.

\subsection{Fűtési igény}

\begin{lstlisting}
GET /heatingRequest?name=PadloKozos&setTemp=true
GET /heatingRequest?name=PadloKozos&setTemp=false
\end{lstlisting}

A \texttt{name} konfigurált fűtési elemnév, a \texttt{setTemp} ebben az API-ban logikai fűtési igényt jelent.

\subsection{Manuális pumpa- és IO-vezérlés}

\begin{lstlisting}
POST /manualControlPumps
Content-Type: application/json
\end{lstlisting}

Aktuális formátum:

\begin{lstlisting}[language=json]
{
  "3": {"manual": true, "value": 50},
  "4": {"manual": false}
}
\end{lstlisting}

Kompatibilitási formátum is támogatott:

\begin{lstlisting}[language=json]
{"3": 50, "4": 0}
\end{lstlisting}

A PWM értéknek 0 és 100 közé kell esnie. A manuális mód felülírja az automatikus vezérlést az adott IO-n.

\subsection{Eszközinformáció}

\begin{lstlisting}
GET /deviceInfo
\end{lstlisting}

A válasz tartalmazza a firmware uptime értékét, WiFi adatokat, chip- és flash-információkat, SPIFFS teljes/használt/szabad méretét, heap adatokat és reset információt.

\subsection{Konfiguráció lekérése és feltöltése}

\begin{lstlisting}
GET /configuration
POST /configuration
\end{lstlisting}

A GET a jelenlegi \texttt{/config.json} fájlt adja vissza. A POST multipart file upload, a fájl neve legyen \texttt{config.json}. A firmware validálja a konfigurációt, biztonsági mentést készít, majd a sikeres feltöltés után újraindítás szükséges.

\subsection{Újraindítás}

\begin{lstlisting}
GET /reboot?user=admin&password=password
\end{lstlisting}

Ez a jelenlegi implementációban rögzített hitelesítést használ. Éles rendszerben ezt módosítani kell.

\section{Bináris állapotnapló}

A naplózás a \texttt{StateDataLogger} osztályban van. A fájlnév:

\begin{lstlisting}
/state_YYYY-MM-DD.bin
\end{lstlisting}

A napló az NTP kliens által szolgáltatott dátum alapján napi fájlba kerül. A naplózás legalább másodpercenként ellenőriz, de csak változás esetén ír új mintát. A hőmérsékleti küszöb 0,1~$^\circ$C, build flaggel módosítható.

\subsection{Fájlfejléc}

\begin{longtable}{rll}
\toprule
Offset & Méret & Tartalom \\
\midrule
0 & 4 bájt & \texttt{HLOG} magic \\
4 & 1 bájt & formátumverzió, jelenleg 1 \\
5 & 3 bájt & fenntartott \\
\bottomrule
\end{longtable}

\subsection{Rekordfejléc}

Minden rekordhoz 5 bájtos fejléc tartozik:

\begin{longtable}{rll}
\toprule
Méret & Típus & Jelentés \\
\midrule
1 bájt & \texttt{uint8\_t} & rekordtípus: 1 séma, 2 minta \\
2 bájt & \texttt{uint16\_t} & payload hossza, little-endian \\
2 bájt & \texttt{uint16\_t} & CRC16 a payloadra \\
\bottomrule
\end{longtable}

\subsection{Sémarekord}

A séma azonosítókat rendel az emberileg értelmezhető kulcsokhoz. Egy mező leírása 2 bájt ID-t, 1 bájt típust, 1 bájt névhosszt és a mezőnév bájtjait tartalmazza. Példák:

\begin{lstlisting}
io/3
sensor/2814576B00000000059
state/Puffer950l/state
state/Puffer950l/HT
\end{lstlisting}

Sémarekord akkor kerül a fájlba, amikor a logger először indul, új nap kezdődik, vagy az aktuális mezőkészlet/típus megváltozik. A konfigurációból később megjelenő új szenzor új sémát eredményezhet.

\subsection{Mintarekord}

Egy mintarekord egy időpontot és az adott időpontban megváltozott mezőket tartalmazza. Az első időbélyeg teljes 32 bites Unix idő, a továbbiak Delta-of-Delta ZigZag-varint formában vannak. Egy mező 2 bájtos azonosítóból, 1 bájtos típusból és 1 vagy 4 bájtos értékből áll.

\begin{itemize}
  \item \texttt{uint8}: IO/PWM vagy enum/állapotérték;
  \item \texttt{float32}: hőmérséklet és \texttt{HT} érték;
  \item egy mintában több mező osztozik az időbélyegen és a CRC-n.
\end{itemize}

\subsection{JSON visszaolvasás}

\begin{lstlisting}
GET /state-history/files
GET /state-history?date=YYYY-MM-DD
DELETE /state-history?date=YYYY-MM-DD
\end{lstlisting}

A fájllista-végpont az eszköz aktuális NTP-dátumát és a megtalált naplófájlokat adja vissza:

\begin{lstlisting}[language=json]
{
  "currentDate": "2026-11-06",
  "files": [
    {"name": "state_2026-11-05.bin", "size": 6213}
  ]
}
\end{lstlisting}

A frontendnek a \texttt{currentDate} mezőt kell használnia alapértelmezett dátumként. A lekérdezés válasza változáslistát ad:

\begin{lstlisting}[language=json]
[
  {
    "timestamp": 1793913984,
    "changes": {
      "io/1": 0,
      "sensor/2814576B00000000059": 48.5,
      "state/Puffer950l/state": "CHARGE"
    }
  }
]
\end{lstlisting}

A frontendnek nem szabad a \texttt{changes} objektumot teljes állapotként kezelnie. A rekordokat időrendben kell feldolgozni, és minden változást az előző teljes állapotba kell merge-elni.

\section{SPIFFS és tárhely}

ESP8266 esetén a \texttt{esp12e} környezet alapértelmezett flash-elrendezést használ. ESP32 esetén a \texttt{partitions\_4mb\_web.csv} SPIFFS partíciója \texttt{0x230000} bájt, azaz 2\,293\,760 bájt. A ténylegesen szabad helyet a \texttt{GET /deviceInfo} válasz \texttt{spiffs.total}, \texttt{spiffs.used} és \texttt{spiffs.free} mezői adják.

A naplófájlok automatikus törlése jelenleg nincs bekapcsolva. Több nap után a frontend vagy egy karbantartási feladat használhatja a DELETE végpontot.

\subsection{Szöveges alkalmazáslog}

A \texttt{LOG\_WRITE} makró a soros kimenet mellett az aktuális napi szöveges fájlba is ír. A fájlok neve:

\begin{lstlisting}
/log_YYYY-MM-DD.txt
\end{lstlisting}

A \texttt{/log.txt} régi/kompatibilitási fájl is beleszámít a kvótába. Az összes \texttt{log*.txt} fájl együttes maximális mérete 10\,240 bájt. Minden új logbejegyzés előtt a rendszer összegyűjti a logfájlok méretét, majd név szerinti sorrendben törli a legrégebbi fájlokat, amíg az új bejegyzés be nem fér.

Ez a kvóta csak a szöveges alkalmazáslogokra vonatkozik. A \texttt{state\_*.bin} állapotnaplók külön fájlok, és nem számítanak bele ebbe a 10 KB-os limitbe.

A fájlműveletek csökkentése érdekében a logüzenetek először egy 2048 bájtos RAM pufferbe kerülnek. A puffer kiürül:

\begin{itemize}
  \item ha megtelik;
  \item 5 másodperccel az első pufferelt üzenet után;
  \item dátumváltáskor;
  \item közvetlen, kényszerített flush esetén.
\end{itemize}

A soros kimenet ettől függetlenül azonnali. A pufferméret és az időzített flush build flaggel módosítható a \texttt{APPLICATION\_LOG\_BUFFER\_SIZE} és \texttt{APPLICATION\_LOG\_FLUSH\_INTERVAL\_MS} makrókkal.

\section{Build és telepítés}

Gyakori környezetek a \texttt{platformio.ini} fájlban:

\begin{longtable}{p{0.35\textwidth}p{0.52\textwidth}}
\toprule
Környezet & Cél \\
\midrule
\texttt{esp12e} & ESP8266 normál firmware \\
\texttt{esp12e\_testOverAPI} & ESP8266 API integrációs tesztekhez \\
\texttt{esp12e\_testRandomData} & ESP8266 szintetikus/random adatokhoz \\
\texttt{esp32dev} & ESP32 normál firmware \\
\texttt{esp32dev\_testOverAPI} & ESP32 API integrációs tesztekhez \\
\texttt{linux\_test} & natív unit teszt környezet \\
\bottomrule
\end{longtable}

\begin{lstlisting}
pio run -e esp12e_testOverAPI
pio run -e esp12e_testOverAPI -t upload
pio run -e esp12e_testOverAPI -t uploadfs
\end{lstlisting}

A naplózási build flag-ek:

\begin{lstlisting}
-DSTATE_LOG_TEMP_THRESHOLD_C=0.1f
-DSTATE_LOG_MIN_CAPTURE_INTERVAL_MS=1000UL
-DAPI_MIN_REQUEST_INTERVAL_MS=0UL
\end{lstlisting}

\section{Tesztelés}

A Python integrációs tesztek a \texttt{tests/} könyvtárban vannak. A firmware-nek hálózaton elérhetőnek kell lennie.

\begin{lstlisting}
python3 -m venv .venv
source .venv/bin/activate
python -m pip install -r requirements-test.txt
HEATING_CONTROLLER_URL=http://192.168.0.106 \\
  python -m pytest -m integration -vv -s
\end{lstlisting}

A tesztek alapértelmezett beállási várakozása 3 másodperc. A részletes API riportokhoz használható a \texttt{--api-trace}, \texttt{--api-report} és \texttt{--api-html-report} kapcsoló.

\section{Továbbfejlesztési javaslatok}

\begin{enumerate}
  \item Induláskor a meglévő fájl utolsó sémájának beolvasása, hogy újraindításkor ne kerüljön be feleslegesen azonos sémarekord.
  \item A fájllista végpont kiegészítése teljes, használt és szabad SPIFFS mérettel.
  \item Napló-retenció vagy maximális naplófájl-szám bevezetése.
  \item A \texttt{/reboot} hitelesítésének konfigurálhatóvá tétele.
  \item A state enumok központi, verziózott dictionary-je az ismeretlen állapotok hash-elése helyett.
  \item A bináris formátumhoz host oldali dekóder és automatikus formátumteszt készítése.
\end{enumerate}

\end{document}
